\section{What is left}\label{sec:reduction}

This section assembles the obstructions into a statement about where the Hodge
conjecture can still be attacked. The statement is an equivalence and not a
reduction, in the sense that neither side is easier than the other; its value
is that it names the single place where a new idea has to enter, and
\Cref{thm:hyperplane} says that the classical inductive handle is blind
exactly there.

\subsection{Reduction to primitive classes}

\begin{proposition}[Lefschetz decomposition of Hodge classes]%
\label{prop:lefschetzdecomp}
Let $X$ be smooth projective of dimension $d$ with ample class $L$. Every
$\alpha\in\Hdg^{p}(X)$ can be written
\[
  \alpha=\sum_{r\ge\max(0,\,p-d+p)}L^{r}\beta_{r},
  \qquad
  \beta_{r}\in\Hdg^{p-r}(X)\cap H^{2p-2r}_{\mathrm{prim}}(X,\QQ),
\]
and the $\beta_{r}$ are uniquely determined. Consequently the Hodge conjecture
holds for $X$ if and only if every primitive rational Hodge class on $X$ is
algebraic.\footnote{The rationality of $L$ is what makes the primitive
parts rational. For a real K\"ahler class that is not a multiple of a rational
one the primitive subspaces are real but not rational, and a rational Hodge
class need not decompose into rational primitive pieces.}
\end{proposition}

\begin{proof}
The Lefschetz decomposition
$H^{k}(X,\QQ)=\bigoplus_{r}L^{r}H^{k-2r}_{\mathrm{prim}}(X,\QQ)$ is a
decomposition of rational Hodge structures, because $L$ is of type $(1,1)$ and
the primitive parts are kernels of morphisms of Hodge structures. A Hodge
class therefore has Hodge classes as its components. For the last statement,
if each primitive $\beta_{r}$ is algebraic then so is $L^{r}\beta_{r}$, since
$L$ is algebraic and algebraic classes form a subring, and hence so is
$\alpha$; the converse is trivial.
\end{proof}

\subsection{The equivalence}

\begin{theorem}\label{thm:equivalence}
The following are equivalent.
\begin{enumerate}[label=\textup{(\roman*)}]
\item The Hodge conjecture holds for every smooth complex projective variety.
\item The Lefschetz standard conjecture $B(X)$ holds for every smooth complex
  projective variety $X$, and every primitive rational Hodge class of middle
  degree on every smooth complex projective variety of even dimension is
  algebraic.
\end{enumerate}
\end{theorem}

\begin{proof}
(i) implies (ii). The second half of (ii) is a special case of (i). For
$B(X)$, the Hodge conjecture for $X\times X$ implies that the K\"unneth
components of the class of the diagonal that realise $\Lambda$ are algebraic:
$\Lambda$ is a morphism of Hodge structures of bidegree $(-1,-1)$, hence is
given by a Hodge class on $X\times X$ by the standard correspondence between
morphisms of Hodge structures and Hodge classes on the product, and (i)
applied to $X\times X$ makes that class algebraic. This is Kleiman's argument
\cite[\S2]{Kle68}.

(ii) implies (i). We argue by induction on $d=\dim X$. For $d\le1$ the Hodge
conjecture is classical. Let $d\ge2$ and assume the conjecture for all smooth
projective varieties of dimension less than $d$. By
\Cref{prop:lefschetzdecomp} it suffices to prove that a primitive Hodge class
$\alpha\in H^{2p}(X,\QQ)$ is algebraic, and by hard Lefschetz we may assume
$2p\le d$.

If $2p=d$, then $\alpha$ is a primitive Hodge class of middle degree on a
variety of even dimension, so it is algebraic by the second half of (ii).

If $2p<d$, choose a smooth hypersurface $Y\in|mH|$ for an ample $H$ and
$m\gg0$, which exists by Bertini, and put $L=c_{1}(\cO_{X}(Y))$. Since
$2p\le d-2$, \Cref{thm:LHT} makes $j^{*}\colon H^{2p}(X,\QQ)\to H^{2p}(Y,\QQ)$
an isomorphism, and $j^{*}\alpha$ is a rational Hodge class on $Y$ because
$j^{*}$ is a morphism of Hodge structures. As $\dim Y=d-1<d$, the inductive
hypothesis makes $j^{*}\alpha$ algebraic, say $j^{*}\alpha=\cl(W)$. Then
$\cl(j_{*}W)=j_{*}j^{*}\alpha=L\alpha$ by \Cref{lem:gysinrestrict}, so
$L\alpha$ is algebraic. By $B(X)$ the operator $\Lambda$ carries algebraic
classes to algebraic classes, and by \Cref{lem:LambdaL},
$\alpha=(d-2p)^{-1}\Lambda(L\alpha)$, which is therefore algebraic. Note
$d-2p\ge2$, so the scalar is nonzero.
\end{proof}

\begin{remark}[This is not a simplification]\label{rem:notsimpler}
Both halves of (ii) are open, and $B(X)$ is itself a consequence of (i), so
\Cref{thm:equivalence} does not reduce the problem to something smaller. What
it does is locate the difficulty. The induction in the proof of (ii) implies
(i) runs smoothly through every degree except the middle one, where it stops
and hands the work to the hypothesis; and \Cref{thm:hyperplane} says that the
step it uses below the middle degree, restriction to a hypersurface, returns
zero at the middle degree, so the induction cannot be extended to cover it by
any refinement of the same step.
\end{remark}

\begin{corollary}[Where a new idea must enter]\label{cor:whereidea}
Any proof of the Hodge conjecture must, at some point, produce an algebraic
cycle representing a primitive Hodge class of middle degree on a variety of
even dimension, by a method that does not factor through restriction to
hyperplane sections of that variety.
\end{corollary}

\begin{proof}
Immediate from \Cref{thm:equivalence} and \Cref{cor:blind}.
\end{proof}

\begin{remark}[Methods that do not factor through hyperplane sections]%
\label{rem:othermethods}
\Cref{cor:whereidea} is not a statement that the conjecture is unapproachable,
and we list the known methods that are outside its scope, since they are where
the successful cases come from. The Kuga--Satake construction replaces a
weight-two Hodge structure of K3 type by the $H^{1}$ of an abelian variety, and
the class it produces is of type $(1,1)$, hence algebraic by Lefschetz; this
is not a restriction to a subvariety and is untouched. Degenerations and
limits of Hodge structures move the variety rather than cutting it.
Correspondences from other varieties, of which the Kuga--Satake construction
is one instance and Markman's construction another, are also outside the
scope. Finally, arguments through the coniveau filtration, which show that a
class supported on a proper closed subvariety is controlled by that
subvariety, are compatible with everything here: they lower the dimension
legitimately, because the class is actually supported on the smaller variety,
which is exactly what a primitive coniveau-zero class is not.
\end{remark}

\begin{figure}[htbp]
\centering
  \includegraphics[width=0.92\textwidth]{fig_moduli}
\caption{The parameter count behind \Cref{thm:equivalence}. For
$n=1,\dots,6$ each prism has the height $n(2n+1)$ of the moduli of principally
polarised abelian $2n$-folds; its lower part, of height $n^{2}$, is the Weil
family of \Cref{prop:weildomain}, and its upper part is the codimension
$n(n+1)$. The gap widens quadratically, so the Weil locus is thin,
and by \Cref{cor:weilalgebraic} it is an algebraic subvariety. Both facts
matter for what is left: the locus one would have to deform along is small and
is a variety, and neither of those things produces a cycle on it.}
\label{fig:roadmap}
\end{figure}
